% -*- TeX:UK -*-
\documentclass[10pt]{beamer}
\usetheme{metropolis}
%\useinnertheme{rectangles}'
\setbeamercovered{%
still covered={\opaqueness<1->{15}},
again covered={\opaqueness<1->{40}}}

\hypersetup{colorlinks,linkcolor=black,urlcolor=brown,citecolor=brown}

\usepackage{wasysym}
\usepackage{amsmath,amssymb,amsthm}
\usepackage{unicode-math}

% We set the Lucida OTF fonts as default
\usepackage{fontspec}
\setmainfont{Lucida Bright OT}
\setsansfont{Lucida Sans OT}
\setmonofont{Lucida Console DK}[Scale=MatchLowercase]

\newfontfamily\typiconsfont{Typicons}[Scale=0.9]
\newcommand\typicons[1]{{\typiconsfont\symbol{#1}}}
\newcommand\Attention{\colorbox{yellow}{\textcolor{black}{\typicons{"E136}}}\xspace}
\newcommand\Advanced{\colorbox{black}{\textcolor{white}{\typicons{"E04E}}}\xspace}

\newfontfamily\webglyphsfont{WebHostingHub-Glyphs}[Scale=0.7]
\newcommand\webglyphs[1]{{\webglyphsfont\symbol{#1}}}
\newcommand\Discussion{\colorbox{white}{\textcolor{black}{\webglyphs{"F134}}}\xspace}
\newcommand\DExamples{\colorbox{black}{\textcolor{white}{\webglyphs{"F134} examples?}}}
\newcommand\Reading{\colorbox{black}{\textcolor{white}{\webglyphs{"F0C1}}}\xspace}

\newfontfamily\meteoconsfont{Meteocons}
\newcommand\meteocons[1]{{\meteoconsfont\symbol{#1}}}
\newcommand\meteosun{\meteocons{"0042}}
\newcommand\meteosuncloud{\meteocons{"0048}}
\newcommand\meteorain{\meteocons{"0052}}
\newcommand\meteosolidsun{\meteocons{"0031}}

\newfontfamily\KRfarmfont{KR Back On The Farm}
\newcommand\KRfarm[1]{{\KRfarmfont\symbol{#1}}}
\newcommand\farmplant{\KRfarm{"0049}}

\usepackage{polyglossia}
\setdefaultlanguage[variant = british, ordinalmonthday = false]{english}

\usepackage[style=authoryear-comp,firstinits,sortcites,maxcitenames=2,%
    mincitenames=1,maxbibnames=10,minbibnames=10,uniquename=mininit,%
    uniquelist=minyear,sortfirstinits=true]{biblatex}
\addbibresource{../references/ecophys.bib}
%\renewcommand{\bibfont}{\small}

\usepackage{abbrev}
\DeclareGraphicsExtensions{.png,.jpg,.jpeg,.eps,.pdf}

\begin{document}

\title{PBIO-141\\Sensory and Physiological Ecology of  Plants}
\subtitle{12b: Carbon Assimilation (E)}
\author{Pedro J. Aphalo}
\date{January-February 2018}
\institute[Univ.\ of Helsinki]{M.Sc.\ in Plant Biology, University of Helsinki\\[2ex] \url{http://blogs.helsinki.fi/aphalo/}}

  \begin{frame}
    \maketitle
  \end{frame}

  \begin{frame}[c]
    \begin{center}
      \begin{small}
        \copyright 2006-2018 by Pedro J. Aphalo\\
       University of Helsinki, Finland.\\
        \textcolor{blue}{\url{http://blogs.helsinki.fi/aphalo/}}\\[2ex]
      \end{small}

      \begin{footnotesize}
        Sensory and Physiological Ecology of  Plants: `12b: Carbon Assimilation (E)' by Pedro J. Aphalo is licensed under a Creative Commons Attribution-ShareAlike 4.0 International License.

        Illustrations and text quoted from copyrighted sources is excluded from this license and their use should respect the original licenses.\\[2ex]
      \end{footnotesize}

      \centering\includegraphics[width=6em]{../figures/copyright/by-sa}
    \end{center}
  \end{frame}

  \begin{frame}
    \frametitle{Outline}
    \tableofcontents
  \end{frame}

\section{Canopy photosynthesis}

\begin{frame}{Canopy photosynthesis}
    \begin{itemize}
        \item If we use the big leaf analogy to express
        photosynthesis of a canopy we can compare the light
        response of a leaf to that of a canopy.%
        \item In a canopy even when upper leaves are light saturated
        (photosynthesizing at maximum rate ($A_{max}$) other lower
        leaves can still increase $A$ if irradiance above the canopy
        increases.
        \item Consequently canopy photosynthesis ($\mathcal{A}$)
        saturates at higher irradiances that $A$ of individual
        leaves.
    \end{itemize}
\end{frame}

\begin{frame}{Needle, shoot and canopy}
    \centering
    \includegraphics[height=0.75\textheight]{figures/Taiz9.11LightResponseNeedleShootCanopy.eps}\\
    {\small The shape of the light response and light saturation
    depends on the scale at which we study the forest
    \autocite[from][]{TaiZei2006}.}
\end{frame}

\begin{frame}{LAI}
    \begin{itemize}
        \item Canopy photosynthesis not only depends on leaf
        photosynthesis.
        \item It depends on leaf area index (LAI), up to a limit the higher the LAI,
        the more light is absorbed.
        \item It depends also on how the leaves are displayed:
        more or less horizontal or vertical.
        \item LAI will depend on environmental conditions.
    \end{itemize}
\end{frame}

\begin{frame}{Soil water affects LAI}
    \centering
    \includegraphics[width=\textwidth]{figures/Crawley1.21LAISoilWater.eps}\\
    {\small Leaf area index (LAI) and site water balance \autocite[from][]{Crawley1997}.}
\end{frame}

\section{Gas-exchange measurement}

\begin{frame}{Gas exchange}
    \begin{itemize}
        \item ``Gas exchange between a leaf and the atmosphere'' is
        the name used to refer to the fluxes of \COtwo and water
        vapour between a leaf and the air surrounding it.
        \item Almost always the flux of water vapour is outwards
        from the leaf and towards the atmosphere. We call it
        \emph{transpiration}.
        \item The flux of \COtwo is in C3 and C4 plants from the
        atmosphere into the leaf during daytime
        (\emph{net photosynthesis}) and from the leaf towards the
        atmosphere (\emph{respiration}) at night. Generic names encompasing them both are:
        \COtwo flux, assimilation, and net carbon exchange (negative for respiration).
    \end{itemize}
\end{frame}

\begin{frame}{Gas-exchange chamber}
    \centering
    \includegraphics[height=0.6\textheight]{figures/leaf_in_chamber.eps}\\
    {\small A leaf or shoot is enclosed in a chamber.\\
    Gas-exchange is now between air in chamber and leaf.}
\end{frame}

\begin{frame}{Closed gas-exchange system}
    \begin{footnotesize}
        \begin{center}
            \fbox{\parbox{3.1cm}{\sf%
            LEAF CHAMBER}}%
            $\makebox[3cm]{\rightarrowfill} \atop
            \makebox[3cm]{%
            \leftarrowfill}$
            \fbox{\parbox{2.2cm}{\sf%
            MEASURING\\ EQUIPMENT}}%
        \end{center}
    \end{footnotesize}
    \begin{itemize}
        \item Some simple gas-exchange systems used for measuring photosynthesis
        are closed systems were there is no air flow through the chamber
        where the leaf in enclosed, except for a loop though instrumentation.

        \item The duration of the measuring period and the differences in concentration
        of \COtwo and water vapour in the chamber air between the beginning and end of
        this measuring period are used to calculate the flows.

        \item From the flows of \COtwo and water and the enclosed leaf area the
        mean \COtwo and water fluxes per unit leaf area can be calculated.
    \end{itemize}
\end{frame}

\begin{frame}{Open gas-exchange system}
\begin{footnotesize}
\begin{center}
\makebox[2cm]{{\sf%
air, \COtwo}%
\rightarrowfill}
\fbox{\parbox{2.3cm}{\sf%
AIR\\ CONDITIONING}}%
$\makebox[3.1cm]{\rightarrowfill} \atop
\makebox[3.1cm]{%
\rightarrowfill
\fbox{\parbox{2.1cm}{\sf%
LEAF CHAMBER}}%
\rightarrowfill}$
\fbox{\parbox{1.8cm}{\sf%
MEASURING\\ EQUIPMENT}}%
\end{center}
\end{footnotesize}
    \begin{itemize}
        \item Most current gas-exchange systems used for measuring photosynthesis
        are open systems in which air flows at a constant rate through a chamber
        where a leaf in enclosed.

        \item The flow rates and the differences in concentration of \COtwo and
        water vapour between the air entering the chamber and that coming
        out of the chamber are used to calculate the flows in or out of the leaf.

        \item From the flows of \COtwo and water in or out of the leaf, and the enclosed leaf area, the
        \COtwo and water fluxes per unit leaf area can be calculated.
    \end{itemize}
\end{frame}

\begin{frame}{Open vs.\ closed systems}
    \begin{description}
        \item[Open system:] Environmental conditions in the chamber do not
        vary during measurement.
        \item[Closed system:] Environmental conditions in the chamber change during measurement.
        With transpiration humidity increases, with photosynthesis \COtwo
        concentration decreases.
    \end{description}
\end{frame}

\begin{frame}{Units}
    \begin{itemize}
        \item Modern: using moles
            \begin{description}
                \item[\COtwo] \umol
                \item[Water vapour] \mmol
            \end{description}
        \item Old: mass
            \begin{description}
                \item[\COtwo] mg cm$^{-2}$ min$^{-1}$ etc.
                \item[Water vapour] g cm$^{-2}$ min$^{-1}$ etc.
            \end{description}
    \end{itemize}
\end{frame}

\begin{frame}{Gas-exchange system}
    \centering
    \includegraphics[height=0.7\textheight]{figures/WalzGFS-3000.eps}\\
    {\small GFS-3000 (Walz, Germany) (\url{http://www.walz.com/}).}
\end{frame}

\begin{frame}{Measuring}
    \centering
    \includegraphics[height=0.7\textheight]{figures/CFS-3000_1big.eps}\\
    {\small from \url{http://www.walz.com/}.}
\end{frame}

\begin{frame}{Measuring}
    \centering
    \includegraphics[height=0.7\textheight]{figures/CFS-3000_6big.eps}\\
    {\small from \url{http://www.walz.com/}.}
\end{frame}

\begin{frame}{Chamber}
    \centering
    \includegraphics[height=0.7\textheight]{figures/GFS-3000_3big.eps}\\
    {\small from \url{http://www.walz.com/}.}
\end{frame}

\begin{frame}{Chamber (measuring)}
    \centering
    \includegraphics[height=0.7\textheight]{figures/CFS-3000_4big.eps}\\
    {\small from \url{http://www.walz.com/}.}
\end{frame}

\begin{frame}{Standard and conifer chambers}
    \centering
    \includegraphics[height=0.5\textheight]{figures/standard.eps}%
    \includegraphics[height=0.5\textheight]{figures/conifer.eps}\\
    {\small Standard chamber (for flat leaves) and conifer chamber (for twigs)
     for the LI-6400 gas-exchange system\\
     (Li-Cor, USA) (from \url{http://www.licor.com/}).}
\end{frame}

\section{Chlorophyll fluorescence}

\begin{frame}{Fluorescence}
    \begin{itemize}
        \item When PAR is absorbed by chlorophyll in a leaf some of
        it is re-emitted as light of a longer wavelength.
        \item This is called fluorescence, and in the case of
        chlorophyll in a leaf fluorescence occurs in the range 680
        nm to 760 nm.
        \item This is the same phenomenon as that on which
        fluorescent tubes are based. The basic emission of mercury
        vapour occurs in the ultraviolet. This radiation is absorbed
        by the coating of the lamp glass tube, and this coating
        fluoresces (at longer wavelengths than the UV)
        and we see this white light.
    \end{itemize}
\end{frame}

\begin{frame}{Fluorescence}
    \begin{itemize}
        \item Not all the energy absorbed by chlorophyll is
        emitted as fluorescence, some is dissipated as heat, as some
        is used in the photochemical reactions of photosynthesis.
        \item We can measure fluorescence under different
        conditions.
        \item At room temperature, in intact leaves most of the
        fluorescence comes from chlorophyll $a$ in PS II.
        \item Consequently by measuring fluorescence at room
        temperature, even in the field we can gather information on
        the functioning of photosynthesis.
    \end{itemize}
\end{frame}

\begin{frame}{Quantum yield}
    \begin{itemize}
        \item In a dark ``adapted'' leaf most PS II reaction centres are
        open (ready to receive excitation), so when measuring
        fluorescence with low intensity light, fluorescence is at a
        minimum ($F_0$).
        \item Under this conditions if we give a high intensity
        light pulse fluorescence is at its maximum ($F_m$), because reaction
        centres quickly close as the electron transport chain is not
        very active.
        \item From these two values the maximum quantum yield of PS II
        ($\phi_{IIm}$) can be calculated as:
        $$\phi_{IIm} = \frac{F_m - F_0}{F_m} = \frac{F_v}{F_m}$$
    \end{itemize}
\end{frame}

\begin{frame}{Chlorophyll fluorescence parameters}
    \centering
    \includegraphics[height=0.7\textheight]{figures/LambersBox4.1Fluorescence.eps}\\
    {\small \autocite[from][]{LambersEtAl1998}.}
\end{frame}

\begin{frame}{$F_v/F_m$ and quantum yield}
    \centering
    \includegraphics[width=\textwidth]{figures/LambersBox4.2FvFmQYield.eps}\\
    {\small \autocite[from][]{LambersEtAl1998}.}
\end{frame}

\begin{frame}{In the light}
    \begin{itemize}
        \item A quantum yield of PS II in the light ($\phi_{II}$)
        can be calculated from fluorescence measured in light
        ``adapted'' (some minutes) leaves as:
        $$\phi_{II} = \frac{\Delta F}{F_m^\prime}$$
        \item In this case $\phi_{II} < \phi_{IIm}$.
        \item The electron transport rate ($J_F$) can also be calculated, if
        the absorbed radiation is available (e.g. calculated from irradiance
        and leaf absorptance).
    \end{itemize}
\end{frame}

\begin{frame}{Fluorescence and assimilation}
    \centering
    \includegraphics[height=0.75\textheight]{figures/LambersBox4.3FluoAssim.eps}\\
    {\small \autocite[from][]{LambersEtAl1998}.}
\end{frame}

\begin{frame}{Fluorescence results}
    \centering
    \includegraphics[height=0.7\textheight]{figures/Huang2006Fig3LightUtilizationSoybean.eps}\\
    {\small Light utilization on a sunny day \autocite[from][]{HuangEtAl2006}.}
\end{frame}

\begin{frame}{Uses}
    \begin{itemize}
        \item The biophysical details of fluorescence in leaves
        are not well understood yet.
        \item However, fluorescence based measurements are very
        useful as they are non-destructive and can, at least at
        ambient temperature, be done in the field with portable
        instruments.
        \item More sophisticated and bulky instruments can be
        used in the laboratory.
    \end{itemize}
\end{frame}

\section{Measuring fluorescence}

\begin{frame}{MINI-PAM}
    \centering
    \includegraphics[height=0.65\textheight]{figures/WalzMINI-PAM.eps}\\
    {\small Instrument to measure chlorophyll fluorescence parameters
    in the field (Walz, Germany) (from \url{http://www.walz.com/}).}
\end{frame}

\begin{frame}{Imaging}
    \begin{itemize}
        \item A more recent development are instruments for
        imaging fluorescence measurements and the parameters derived
        from these measurements (in real time).
        \item For example we may get an image of a leaf where
        different colours indicate the value of $F_v/F_m$ in
        different parts of the leaf.
        \item Such imaging is available from whole leaves to
        tissues observed under the microscope.
    \end{itemize}
\end{frame}

\begin{frame}{Imaging of Chl fluorescence params.}
    \centering
    \includegraphics[height=0.7\textheight]{figures/WalzImagingPAM-2.eps}\\
    {\small  from \url{http://www.walz.com/}.}
\end{frame}

\begin{frame}{\ldots also at macro and micro scales}
    \centering
    \includegraphics[height=0.7\textheight]{figures/WalzImagingPAM-1.eps}\\
    {\small  from \url{http://www.walz.com/}.}
\end{frame}

\begin{frame}{Mix and match}
    \begin{itemize}
        \item Nowadays it is possible to mount a fiber optic probe
        on the window of a leaf chamber of a gas-exchange system,
        and measure simultaneously fluorescence and gas-exchange of
        a leaf.
        \item This allows a better description of the state of
        the photosynthetic ``machinery'' in a leaf.
        \item For example for diagnosing photoinhibition in a
        shade plant suddenly exposed to full sunlight, or
        limitations to assimilation in a plant under drought.
    \end{itemize}
\end{frame}

\begin{frame}{Gas-exchange chamber + fluorometer}
    \centering
    \includegraphics[height=0.7\textheight]{figures/WalzGFS-3000fluo.eps}\\
    {\small from \url{http://www.walz.com/}.}
\end{frame}

  \section*{References}
  \begin{frame}[t,allowframebreaks]
    \frametitle{References}
    \printbibliography
  \end{frame}

\end{document}

